% !PRINTED 53
\weberpage{53}
%<<CONT

\[
\alpha_1 = \frac{\eta_1 a_1}{2},\ \alpha_2 = \frac{\eta_2 a_2}{2} \cdots \alpha_\lambda = \frac{\eta_\lambda a_\lambda}{2},\ \alpha_{\lambda+1}=0 \ldots \alpha_\nu = 0,
\]

worin $\eta_1, \eta_2 \ldots \eta_\lambda$ gleich 0 oder gleich 1 sein können, und man erhält alle zweiseitigen Elemente in der Form

\[
A_1^{\frac{\eta_1 a_1}{2}} A_2^{\frac{\eta_2 a_2}{2}} \ldots A_\lambda^{\frac{\eta_\lambda a_\lambda}{2}},
\]

und ihre Anzahl ist, da alle Combinationen von 0 und 1 für die Exponenten $\eta$ zulässig sind, $= 2^\lambda$.

Ist $\chi$ ein beliebiger Charakter von $S$, so ist, wenn $A$ ein zweiseitiges Element ist, $\chi(A) = \pm 1$, da

\[
\chi(A)^2 = \chi(1) = 1
\]

ist.

Ebenso nennen wir einen zweiseitigen Charakter einen solchen, der in der Gruppe der Charaktere ein zweiseitiges Element ist. Da die Gruppe der Charaktere mit der Gruppe $S$ isomorph ist, so giebt es ebenso viele zweiseitige Charaktere als zweiseitige Elemente, nämlich $2^\lambda$. Sie werden nach \S.\,11, (17) dargestellt durch:

\[
\chi_1^{\frac{\eta_1 a_1}{2}} \chi_2^{\frac{\eta_2 a_2}{2}} \ldots \chi_\lambda^{\frac{\eta_\lambda a_\lambda}{2}}.
\]

\gesperrt{Die zweiseitigen Charaktere haben für jedes Element den Werth $\pm 1$.}

Die zweiseitigen Elemente bilden für sich eine Gruppe $T$ vom Grade $2^\lambda$. Ebenso bilden die zweiseitigen Charaktere eine damit isomorphe Gruppe, und man erhält die zu $T$ reciproke Gruppe $U$, wenn man alle Elemente $A$ aufsucht, für die alle zweiseitigen Charaktere den Werth $+1$ haben. Diese Gruppe, die nach \S.\,12, 7. ein Theiler von $S$ vom Index $2^\lambda$ ist, ist nach der letzten Definition weder von der Wahl der Basis noch von den die Charaktere darstellenden Einheitswurzeln abhängig, und ist durch die Natur der Gruppe $S$ vollständig bestimmt. Bezeichnen wir sie mit $G$, so ist das System der Nebengruppen

\weq{3}{G,\ G_1,\ G_2 \ldots G_{2^\lambda - 1}}

dadurch charakterisirt, dass für alle Elemente eines dieser Systeme die zweiseitigen Charaktere ein und dasselbe Werthsystem ergeben. Die Systeme $G, G_1, G_2 \ldots$ werden die in $S$ enthaltenen \gesperrt{Geschlechter} (Genera) genannt. Die Gruppe $G$ speciell heisst das \gesperrt{Hauptgeschlecht}.
